% ==============================================================
% Dekonstruktion des Hard Problems (DHP)
% Im Rahmen der Methodologischen Konventionen (MKO) des OP
% Version 1.2
% Kompiliert mit: pdfLaTeX
% ==============================================================

\documentclass[11pt,a4paper]{article}

% ---------- Grundpakete ----------
\usepackage[utf8]{inputenc}
\usepackage[T1]{fontenc}
\usepackage[german]{babel}
\usepackage{amsmath,amssymb,amsthm}
\usepackage{geometry}
\usepackage{parskip}
\usepackage{enumitem}
\usepackage{booktabs}
\usepackage{xcolor}
\usepackage{hyperref}
\usepackage{csquotes}

% ---------- Layout ----------
\geometry{margin=2.5cm}
\setlist{itemsep=0.2em, topsep=0.2em}
\hypersetup{
    colorlinks=true,
    linkcolor=blue,
    urlcolor=blue,
    pdftitle={Dekonstruktion des Hard Problems (DHP) -- Version 1.2},
    pdfauthor={Forschungsprogramm Ontoph\"{a}nomenalismus}
}

% ---------- Theorem-Umgebungen ----------
\theoremstyle{definition}
\newtheorem{konvention}{Konvention}[section]
\newtheorem{remark}{Bemerkung}[section]
\newtheorem{methodennote}{Methodennote}[section]

\theoremstyle{plain}
\newtheorem{axiom}{Axiom}[section]
\newtheorem{proposition}{Proposition}[section]
\newtheorem{hypothesis}{Hypothese}[section]
\newtheorem{definition}{Definition}[section]

% ---------- Makros f\"{u}r MKO-Kennzeichnungen ----------
\newcommand{\axiomR}[2]{\axiom{\textbf{#1}: #2}}
\newcommand{\axiomP}[2]{\axiom{\textbf{#1}: #2}}
\newcommand{\defi}[2]{\definition{\textbf{#1}: #2}}
\newcommand{\analogie}[2]{\methodennote{\textbf{#1}: #2}}
\newcommand{\hyp}[2]{\hypothesis{\textbf{#1}: #2}}
\newcommand{\bem}[1]{\remark{#1}}
\newcommand{\prop}[2]{\proposition{\textbf{#1}: #2}}

% ---------- Titel ----------
\title{
  \textbf{Dekonstruktion des Hard Problems (DHP)}\\[0.4em]
  {\large Im Rahmen der Methodologischen Konventionen (MKO) des OP}\\[0.5em]
  {\normalsize Version~1.2}
}
\author{
  \textbf{Forschungsprogramm Ontoph\"{a}nomenalismus}\\
  \small Siehe TRANSPARENCY f\"{u}r methodologische Urheberschaft und KI-Einsatz.
}
\date{Juli 2026}

\begin{document}
\maketitle

\begin{abstract}
Dieses Dokument wendet die f\"{u}nf-Ebenen-Sprachregel der Methodologischen Konventionen
(MKO) an, um das \emph{Hard Problem of Consciousness} (Chalmers) zu dekonstruieren.

Es gelten die methodologischen Konventionen gem\"{a}\ss{} MKO. Wir zeigen, dass die
materialistische Zombie-Hypothese eine metaphysische Zusatzannahme darstellt, die weder
empirisch best\"{a}tigt noch logisch erzwungen ist, und dass der Ontoph\"{a}nomenalismus
(OP) die empirisch belegte Korrelation zwischen physikalischer Organisation und
ph\"{a}nomenalem Erleben als ontologischen Ausgangspunkt w\"{a}hlt. Die
Argumentationskette wird mit den MKO-Kennzeichnungen (Axiom R, Axiom P, Definition,
Methodennote, Hypothese) strukturiert, sodass Zirkel-Kontrolle und epistemische
Transparenz explizit werden.

Das Hard Problem erscheint im Rahmen des OP damit nicht mehr als unlösbares R\"{a}tsel
einer einseitigen Emergenz, sondern als ein symmetrisches, projektionsanalytisches Thema.
\end{abstract}

\tableofcontents
\newpage

%==============================================================
\section{Einleitung}
%==============================================================

Das klassische \emph{Hard Problem of Consciousness} (Chalmers) fragt nach der
Erkl\"{a}rung, warum und wie aus einer hochkomplexen physikalischen Struktur ($A$)
subjektives Erleben ($B$) entsteht. In der philosophischen Literatur wird diese Frage
h\"{a}ufig mit der \textbf{philosophischen Zombie-Hypothese} untermauert: Es sei
metaphysisch m\"{o}glich, dass ein System exakt dieselbe physikalische Organisation wie
ein menschliches Gehirn besitzt, jedoch ohne irgendeine Form von Bewusstsein.

Im Rahmen der MKO wollen wir diese Zusatzannahme kritisch pr\"{u}fen und die alternative
Position des OP systematisch darstellen. Eine ausf\"{u}hrliche Bestandsaufnahme aller
L\"{o}sungsversuche und die vollst\"{a}ndige Position des OP finden sich in
\textbf{EG-1 (Der Explanatory Gap)}.

%==============================================================
\section{Die zombische Zusatzannahme}
%==============================================================

\subsection{Materialistische Grundannahme (Axiom R)}

\axiomR{Z1}{Die Setzung der prinzipiellen Trennbarkeit von physikalischer Struktur ($A$)
und ph\"{a}nomenalem Erleben ($B$) im Konstrukt des philosophischen Zombies fungiert im
Materialismus als unhinterfragte, dogmatische Nullhypothese.}
\textit{(Rekonstruktiver, indikativer Schritt)}

\bem{Die Zombie-Hypothese ist ein metaphysisches Postulat, das nicht aus vorherigen
empirischen Pr\"{a}missen folgt, sondern als zus\"{a}tzliche, asymmetrische Annahme
eingef\"{u}hrt wird, um das Privileg der Beweislast abzuw\"{a}lzen.}

\subsection{Empirische Evidenz-L\"{u}cke (Hypothese H)}

\hyp{H1}{Es existiert kein neuro-physiologischer oder empirischer Befund, der ein
funktionsf\"{a}higes, komplex organisiertes Nervensystem ($A$) ohne begleitende
ph\"{a}nomenale Innenperspektive ($B$) nachweist.}

\bem{Die Behauptung ist eine \textbf{Hypothese} (MKO Ebene~5), weil sie empirisch
pr\"{u}fbar und falsifizierbar, im Rahmen unserer gesamten Datenbasis bislang aber
ausnahmslos unbest\"{a}tigt ist.}

\subsection{Methodologische Unbeobachtbarkeit (Methodennote)}

\methodennote{U1}{Subjektives Erleben ist intersubjektiv nicht direkt von au\ss{}en
messbar, sondern zeigt sich physikalisch ausschlie\ss{}lich \"{u}ber seine strukturellen
Korrelate (\textbf{epistemisches Limit}).}

\bem{Dies ist eine \textbf{Methodennote} (MKO Ebene~2), die an den Realismus-Vorbehalt
der MKO erinnert: Alle Aussagen \"{u}ber die ph\"{a}nomenale Innenseite erfordern eine
methodische Konvention zur Interpretation physikalischer Signale.}

\subsection{Logische Konsequenz (Proposition)}

\prop{P1}{Falls die zombische Asymmetrie-Hypothese wahr w\"{a}re, m\"{u}sste die
empirische Kopplung $A \iff B$ eine kontingente, aufhebbare Eigenschaft der Natur sein.
Da dies unbelegt ist, besitzt das Postulat Z1 keinerlei logischen oder empirischen
Vorrang gegen\"{u}ber einer symmetrischen Kopplungstheorie.}

\bem{Die Proposition folgt aus dem Axiom~R (rekonstruktiv) und der Hypothese H1
(empirische Evidenz-L\"{u}cke). Sie ist ein \textbf{logischer Schritt} (MKO Ebene~1),
der das Monopol der materialistischen Nullhypothese bricht.}

%==============================================================
\section{Der Ontoph\"{a}nomenalismus als Gegenposition}
%==============================================================

\subsection{Postulat der ontologischen Grundentscheidung (Axiom P)}

\axiomP{OP1}{Der ausnahmslos beobachtete, bikonditionale Zusammenhang zwischen
physikalischer Struktur ($A$) und Erleben ($B$) wird als unhintergehbarer ontologischer
Ausgangspunkt des OP gesetzt. Jedes theoretische Konstrukt, das $A$ ohne $B$ postuliert,
wird als metaphysisch unbegr\"{u}ndete Zusatzannahme abgewiesen.}

\bem{Dieses Axiom ist ein \textbf{Postulat} (MKO Ebene~3, Typ~P) --- es wird nicht aus
vorherigen materialistischen Schritten abgeleitet, sondern markiert die fundamentale,
sparsame Design-Entscheidung des OP-Systems.}

\subsection{Definition der Korrelation (Definition)}

\defi{Korr}{Die \textbf{Korrelation $K$} bezeichnet das empirisch vorgefundene,
funktionale und bikonditionale Verh\"{a}ltnis zwischen einer komplexen physikalischen
Organisation $A$ und dem ph\"{a}nomenalen Erleben $B$, formal repr\"{a}sentiert durch
das relationale Band $A \iff B$.}

\bem{Definition (MKO Ebene~3) --- sie legt die mathematisch-logische Terminologie fest,
ohne voreilig einseitige Kausalit\"{a}tsbehauptungen zu treffen.}

\subsection{Analogie: Kugel-Schatten-Prinzip (MKO Ebene~4)}

\analogie{A1}{Das ph\"{a}nomenale Erleben ($B$) entspricht einer h\"{o}herdimensionalen
oder metrikfreien informationellen Realit\"{a}t (der \glqq{}Kugel\grqq{}), w\"{a}hrend
die physikalische Struktur ($A$) deren dimensional reduzierte Projektion darstellt (der
\glqq{}Schatten\grqq{}). Die Abbildung $f\colon B \to A$ macht deutlich: Der Schatten
($A$) wird nicht von der Materie erzeugt, sondern er ist die zwingende physikalische
Manifestation der Kugel ($B$).}

\bem{Die Kugel-Schatten-Analogie dient als \textbf{heuristisches Bild} (MKO Ebene~4).
Sie verdeutlicht das erkenntnistheoretische Potenzial: Wer fragt, wie das Gehirn
Bewusstsein erzeugt, k\"{o}nnte --- im Rahmen des OP --- einen Kategorienfehler begehen:
Er fr\"{a}gt, wie der Schatten an der Wand die Kugel im Raum generiert.}

\subsection{Hypothese \"{u}ber empirische Konsequenzen (Hypothese H2)}

\hyp{H2}{In allen physikalisch-biologischen Systemen, in denen die strukturelle
Organisation $A$ nachgewiesen wird, k\"{o}nnte das gleichzeitige Vorhandensein der
ph\"{a}nomenalen Signatur $B$ eine systemische Notwendigkeit darstellen.}

\bem{Dies ist eine im Rahmen zuk\"{u}nftiger KI- und Neurostrukturen \textbf{empirisch
testbare} Aussage (MKO Ebene~5).}

\subsection{Methodennotiz zu Gegenargumenten}

\methodennote{G1}{Klassische materialistische Einw\"{a}nde argumentieren, dass
(i)~hypothetische Epiphänomene oder noch unbekannte physikalische Mechanismen eine
einseitige Kausalit\"{a}t $A \to B$ st\"{u}tzen k\"{o}nnten, oder (ii)~eine reine
Korrelation keine ontologische Identit\"{a}t beweist.}

\bem{Diese Einw\"{a}nde werden als \textbf{Methodennotiz} (MKO Ebene~2) eingeordnet,
da sie den meta-theoretischen Status von Kausalit\"{a}t betreffen, ohne die empirische
Vorrangstellung der Kopplung $K$ zu ersch\"{u}ttern.}

\subsection{Reaktion auf Gegenargumente (Proposition)}

\prop{P2}{Jede alternative reduktionistische Theorie steht in der Pflicht, die
beobachtete Bikonditionalit\"{a}t $K$ ohne Ad-hoc-Zusatzschleifen zu reproduzieren.
Da der OP mit $A \iff B$ die minimale, empirisch deckungsgleiche Struktur w\"{a}hlt,
k\"{o}nnte er gem\"{a}\ss{} Ockhams Rasiermesser erkenntnistheoretisch vorzuziehen sein.}

\bem{Proposition (MKO Ebene~1) --- sie stellt die \textbf{Zirkel-Kontrolle} her, indem
sie zeigt, dass die ph\"{a}nomenologische Position keine unbeobachteten Zombie-Zust\"{a}nde
erfinden muss, die ihr materialistisches Pendant zwingend voraussetzt.}

%==============================================================
\section{Zusammenfassung}
%==============================================================

\begin{enumerate}[label=(\arabic*)]
  \item Die materialistische \textbf{Zombie-Hypothese} ist eine unbegr\"{u}ndete
        metaphysische Zusatzannahme, die weder empirisch gest\"{u}tzt (\textbf{H1})
        noch logisch zwingend ist (\textbf{P1}).
  \item Der \textbf{Ontoph\"{a}nomenalismus} w\"{a}hlt die reale, ungetrennte
        \textbf{Korrelation $K$} als stabilen ontologischen Ausgangspunkt (\textbf{OP1})
        und formalisiert diese Einheit (\textbf{Definition Korr}).
  \item Durch das \textbf{Kugel-Schatten-Prinzip} wird das logische M\"{o}glichkeitsszenario
        einer Umkehrung der Erkl\"{a}rungsrichtung ($B \to A$) anschaulich gemacht, ohne
        es dogmatisch aufzuzwingen (\textbf{A1}). Das Gehirn w\"{a}re in dieser Lesart
        als materielle Signatur des Denkfeldes zu verstehen.
  \item Die empirische Tragf\"{a}higkeit des Modells wird \"{u}ber die
        \textbf{Hypothese H2} an die Realit\"{a}t gekoppelt; jede konkurrierende Theorie
        m\"{u}sste $K$ ohne metaphysischen Ballast reproduzieren (\textbf{P2}).
  \item Die \textbf{Zirkel-Kontrolle} ist explizit: Die Rekonstruktion (R) entlarvt die
        materialistische Voreingenommenheit, w\"{a}hrend das Postulat (P) die
        minimalistische, datentreue Basis des OP fixiert.
\end{enumerate}

Das Hard Problem erscheint damit im OP-Rahmen nicht mehr als unlösbares R\"{a}tsel einer
einseitigen Emergenz, sondern k\"{o}nnte als ein symmetrisches, projektionsanalytisches
Thema verstanden werden, das die methodischen Konventionen der MKO wahrt. Die
vollst\"{a}ndige philosophische Ausarbeitung dieser Position --- einschlie\ss{}lich der
kategorialen Fehlfragen-These, der Zimmer- und Traumanalogie sowie des
Fragedifferentials --- findet sich in \textbf{EG-1 (Der Explanatory Gap)} und
\textbf{DFD (Das Fragedifferential)}.

\end{document}
